% GENERATED FILE. Edit canonical lessons, then run npm run generate:books.
% canonical-source-hash: fnv1a64:b24430b41bb8bc56
% canonical-lessons: FR-C339-preuve, FR-C339-evident, FR-C339-evidence, FR-C339-incontestable, FR-C339-manifeste

\chapter{Les preuves}
\label{ch:fr-les-preuves}
\hlchaptermodality{339}

\begin{chapteropening}
\textbf{By the end of this chapter:} \emph{I can talk about proof: say something is obvious, plain to see, indisputable or clear.}

Complete the last lesson of chapter 339: I can talk about proof: say something is obvious, plain to see, indisputable or clear.
\end{chapteropening}

\section[la preuve]{la preuve — the proof, the evidence}

\label{lesson:FR-C339-preuve}



\begin{quote}
Before the new one: say the French for absurd, then the French for debatable.
\end{quote}

\subsection*{You'll want to know: la preuve}
\textbf{la preuve} — \textquotedblleft{}the proof, the evidence\textquotedblright{}.

\begin{quote}
Tu dis que c'est lui, mais tu n'as aucune preuve.
\end{quote}

\begin{quote}
\emph{You say it was him, but you have no proof.}
\end{quote}

\emph{Faire preuve de} — to show a quality: \emph{faire preuve de courage}.

\begin{cousinweb}
\emph{preuve} is usually traced to Old French \emph{prueve}, a noun formed on \emph{prouver} \textquotedblleft{}to prove\textquotedblright{}, which is inherited from Latin \emph{probāre} \textquotedblleft{}to test, to examine\textquotedblright{} (Lewis \& Short).
\end{cousinweb}

\begin{grammarlens}[title={What you've built}]
One more word for telling, explaining and describing at B1.
\end{grammarlens}

\subsection*{Your turn}
\begin{itemize}
\raggedright
  \item \emph{Say it:} \emph{la preuve}
  \item \emph{Say it:} \emph{la preuve}, once more
  \item \emph{Say it:} a sentence of your own with \emph{la preuve}
  \item \emph{Recall:} say the French for absurd, then the French for debatable, then say \emph{la preuve} again
\end{itemize}

\begin{tcolorbox}[breakable,colback=teal!4,colframe=teal!35!black,title={Before you move on}]
What does la preuve mean? (\textquotedblleft{}The proof, the evidence\textquotedblright{}.) Say it once more.
\end{tcolorbox}
\section[évident]{évident — obvious, evident}

\label{lesson:FR-C339-evident}



\begin{quote}
Before the new one: say the French for debatable, then the French for the proof.
\end{quote}

\subsection*{You'll want to know: évident}
\textbf{évident} — \textquotedblleft{}obvious, evident\textquotedblright{}.

\begin{quote}
C'est évident : il est fatigué.
\end{quote}

\begin{quote}
\emph{It is obvious: he is tired.}
\end{quote}

\emph{Ce n'est pas évident} also means it is not easy.

\begin{cousinweb}
\emph{évident} is a learned borrowing from Latin \emph{ēvidēns} \textquotedblleft{}apparent, visible, clear\textquotedblright{} (EtymDB; Lewis \& Short), built from \emph{ē} \textquotedblleft{}out\textquotedblright{} and \emph{vidēre} \textquotedblleft{}to see\textquotedblright{}.
\end{cousinweb}

\begin{grammarlens}[title={What you've built}]
One more word for telling, explaining and describing at B1.
\end{grammarlens}

\subsection*{Your turn}
\begin{itemize}
\raggedright
  \item \emph{Say it:} \emph{évident}
  \item \emph{Say it:} \emph{évident}, once more
  \item \emph{Say it:} a sentence of your own with \emph{évident}
  \item \emph{Recall:} say the French for debatable, then the French for the proof, then say \emph{évident} again
\end{itemize}

\begin{tcolorbox}[breakable,colback=teal!4,colframe=teal!35!black,title={Before you move on}]
What does évident mean? (\textquotedblleft{}Obvious, evident\textquotedblright{}.) Say it once more.
\end{tcolorbox}
\section[l'évidence]{l'évidence — the obvious, what is plain to see}

\label{lesson:FR-C339-evidence}



\begin{quote}
Before the new one: say the French for the proof, then the French for obvious.
\end{quote}

\subsection*{You'll want to know: l'évidence}
\textbf{l'évidence} — \textquotedblleft{}the obvious, what is plain to see\textquotedblright{}.

\begin{quote}
Il faut se rendre à l'évidence : nous sommes perdus.
\end{quote}

\begin{quote}
\emph{We have to face the facts: we are lost.}
\end{quote}

A false friend: \emph{l'évidence} is what is plain to see; evidence in court is \emph{la preuve}. \emph{Se rendre à l'évidence} — to face the facts.

\begin{cousinweb}
\emph{évidence} is a learned borrowing from Latin \emph{ēvidentia} \textquotedblleft{}clearness, distinctness\textquotedblright{} (Lewis \& Short), from \emph{ēvidēns} \textquotedblleft{}clear\textquotedblright{}, built on \emph{vidēre} \textquotedblleft{}to see\textquotedblright{}.
\end{cousinweb}

\begin{grammarlens}[title={What you've built}]
One more word for telling, explaining and describing at B1.
\end{grammarlens}

\subsection*{Your turn}
\begin{itemize}
\raggedright
  \item \emph{Say it:} \emph{l'évidence}
  \item \emph{Say it:} \emph{l'évidence}, once more
  \item \emph{Say it:} a sentence of your own with \emph{l'évidence}
  \item \emph{Recall:} say the French for the proof, then the French for obvious, then say \emph{l'évidence} again
\end{itemize}

\begin{tcolorbox}[breakable,colback=teal!4,colframe=teal!35!black,title={Before you move on}]
What does l'évidence mean? (\textquotedblleft{}The obvious, what is plain to see\textquotedblright{}.) Say it once more.
\end{tcolorbox}
\section[incontestable]{incontestable — indisputable, unquestionable}

\label{lesson:FR-C339-incontestable}



\begin{quote}
Before the new one: say the French for obvious, then the French for the obvious.
\end{quote}

\subsection*{You'll want to know: incontestable}
\textbf{incontestable} — \textquotedblleft{}indisputable, unquestionable\textquotedblright{}.

\begin{quote}
Son talent est incontestable.
\end{quote}

\begin{quote}
\emph{Her talent is indisputable.}
\end{quote}

Literally that no one can contest. One form for both genders.

\begin{cousinweb}
\emph{incontestable} is built from \emph{in-} \textquotedblleft{}not\textquotedblright{}, \emph{contester} and \emph{-able}. \emph{contester} comes from Latin \emph{contestārī} \textquotedblleft{}to call to witness\textquotedblright{} (Lewis \& Short), built on \emph{testis} \textquotedblleft{}witness\textquotedblright{}.
\end{cousinweb}

\begin{grammarlens}[title={What you've built}]
One more word for telling, explaining and describing at B1.
\end{grammarlens}

\subsection*{Your turn}
\begin{itemize}
\raggedright
  \item \emph{Say it:} \emph{incontestable}
  \item \emph{Say it:} \emph{incontestable}, once more
  \item \emph{Say it:} a sentence of your own with \emph{incontestable}
  \item \emph{Recall:} say the French for obvious, then the French for the obvious, then say \emph{incontestable} again
\end{itemize}

\begin{tcolorbox}[breakable,colback=teal!4,colframe=teal!35!black,title={Before you move on}]
What does incontestable mean? (\textquotedblleft{}Indisputable, unquestionable\textquotedblright{}.) Say it once more.
\end{tcolorbox}
\section[manifeste]{manifeste — clear, obvious, manifest}

\label{lesson:FR-C339-manifeste}



\begin{quote}
Before the new one: say the French for the obvious, then the French for indisputable.
\end{quote}

\subsection*{You'll want to know: manifeste}
\textbf{manifeste} — \textquotedblleft{}clear, obvious, manifest\textquotedblright{}.

\begin{quote}
Il y a une erreur manifeste dans la facture.
\end{quote}

\begin{quote}
\emph{There is a clear mistake in the bill.}
\end{quote}

Formal. As a noun, \emph{un manifeste} is a manifesto.

\begin{cousinweb}
\emph{manifeste} is a learned borrowing from Latin \emph{manifestus} \textquotedblleft{}palpable, plain, evident\textquotedblright{} (Lewis \& Short), which the dictionary links to \emph{manus} \textquotedblleft{}hand\textquotedblright{}: so clear you could touch it.
\end{cousinweb}

\begin{grammarlens}[title={What you've built}]
One more word for telling, explaining and describing at B1.
\end{grammarlens}

\subsection*{Your turn}
\begin{itemize}
\raggedright
  \item \emph{Say it:} \emph{manifeste}
  \item \emph{Say it:} \emph{manifeste}, once more
  \item \emph{Say it:} a sentence of your own with \emph{manifeste}
  \item \emph{Recall:} say the French for the obvious, then the French for indisputable, then say \emph{manifeste} again
\end{itemize}

\begin{tcolorbox}[breakable,colback=teal!4,colframe=teal!35!black,title={Before you move on}]
What does manifeste mean? (\textquotedblleft{}Clear, obvious, manifest\textquotedblright{}.) Say it once more.
\end{tcolorbox}
